\documentclass[12pt,a4paper]{article}

\usepackage[T2A]{fontenc}
\usepackage[utf8]{inputenc}
\usepackage[russian]{babel}
\usepackage{amsmath,amssymb,amsthm,mathtools}
\usepackage{geometry}
\usepackage{microtype}
\usepackage{enumitem}
\usepackage{xcolor}
\usepackage{tikz}
\usepackage{hyperref}

\geometry{left=25mm,right=20mm,top=22mm,bottom=24mm}
\hypersetup{
  colorlinks=true,
  linkcolor=blue!45!black,
  urlcolor=blue!55!black
}
\setlist{nosep}
\setlength{\parindent}{1.25em}
\setlength{\parskip}{0.15em}

\newtheorem{theorem}{Теорема}
\newtheorem{proposition}{Утверждение}
\theoremstyle{definition}
\newtheorem{definition}{Определение}
\theoremstyle{remark}
\newtheorem*{remark}{Замечание}

\DeclareMathOperator{\Sp}{Sp}
\DeclareMathOperator{\diag}{diag}
\DeclareMathOperator{\cond}{cond}
\newcommand{\CC}{\mathbb{C}}
\newcommand{\RR}{\mathbb{R}}
\newcommand{\ZZ}{\mathbb{Z}}
\newcommand{\eps}{\varepsilon}
\newcommand{\norm}[1]{\left\lVert #1\right\rVert}

\title{Циклические матрицы и свёртки на группе $\ZZ_n$}
\author{Конспект лекции}
\date{19 сентября 2026 года}

\begin{document}

\maketitle

\section{Обозначения и связь с предыдущей лекцией}

Работаем в пространстве $\CC^n$. Координаты векторов и индексы матриц
нумеруем от $0$ до $n-1$. Единичную матрицу обозначаем $E=I_n$, а через
$V$ обозначаем матрицу циклического сдвига:
\[
  V(x_0,x_1,\dots,x_{n-1})^{\mathsf T}
  =(x_{n-1},x_0,\dots,x_{n-2})^{\mathsf T},
  \qquad V^n=E.
\]
Напомним, что циклическая матрица с первым столбцом
$a=(a_0,a_1,\dots,a_{n-1})^{\mathsf T}$ имеет вид
\[
  A=\begin{pmatrix}
    a_0 & a_{n-1} & a_{n-2} & \cdots & a_1\\
    a_1 & a_0 & a_{n-1} & \cdots & a_2\\
    a_2 & a_1 & a_0 & \ddots & \vdots\\
    \vdots & \vdots & \ddots & \ddots & a_{n-1}\\
    a_{n-1} & a_{n-2} & \cdots & a_1 & a_0
  \end{pmatrix}.
\]
Каждый следующий столбец получается циклическим сдвигом предыдущего.
Символом матрицы называется многочлен
\[
  a(t)=\sum_{r=0}^{n-1}a_rt^r.
\]
На предыдущей лекции было получено представление
\[
  A=a(V)=a_0E+a_1V+\dots+a_{n-1}V^{n-1}.
\]

Корни $n$-й степени из единицы и соответствующие собственные векторы
матрицы $V$ обозначаем следующим образом:
\[
  \eps_k=e^{2\pi i k/n},\qquad
  v_k=(1,\eps_k^{-1},\dots,\eps_k^{-(n-1)})^{\mathsf T},
  \qquad Vv_k=\eps_kv_k.
\]
Векторы $v_0,\dots,v_{n-1}$ образуют базис. Матрица
\[
  F=(v_0\ v_1\ \cdots\ v_{n-1}),\qquad
  F_{jk}=e^{-2\pi i jk/n},
\]
является матрицей дискретного преобразования Фурье в принятом соглашении
о знаке и нормировке.

\section{Спектр и диагонализация циклической матрицы}

\subsection{Завершение теоремы об отображении спектра}

Для матрицы $B\in\CC^{n\times n}$ и многочлена $p$ теорема об отображении
спектра утверждает, что
\[
  \Sp(p(B))=p(\Sp(B)).
\]
В начале рукописи заканчивается доказательство обратного включения.
Восстановим этот переход.

Пусть $\mu\in\Sp(p(B))$, то есть $p(B)-\mu E$ необратима. Если $p$
непостоянен, над полем $\CC$ можно записать
\[
  p(t)-\mu=c\prod_{r=1}^{m}(t-t_r),\qquad c\ne0.
\]
Тогда
\[
  p(B)-\mu E=c\prod_{r=1}^{m}(B-t_rE).
\]
Если бы все множители были обратимы, их произведение также было бы
обратимо. Поэтому для некоторого $r$ матрица $B-t_rE$ необратима,
откуда $t_r\in\Sp(B)$. При этом $p(t_r)=\mu$, и значит,
$\mu\in p(\Sp(B))$. Для постоянного многочлена равенство очевидно.

\subsection{Собственные числа и собственные векторы}

\begin{theorem}
Пусть $A$ --- циклическая матрица с символом $a(t)$. Тогда
\[
  \lambda_k=a(\eps_k),\qquad k=0,1,\dots,n-1,
\]
--- её собственные числа с учётом кратностей, а $v_k$ --- соответствующие
собственные векторы. В частности,
\[
  \Sp(A)=\{a(\eps_0),a(\eps_1),\dots,a(\eps_{n-1})\}.
\]
\end{theorem}

\begin{proof}
Так как $Vv_k=\eps_kv_k$, для каждой степени выполнено
$V^rv_k=\eps_k^rv_k$. Поэтому
\[
  Av_k=a(V)v_k
  =\sum_{r=0}^{n-1}a_rV^rv_k
  =\left(\sum_{r=0}^{n-1}a_r\eps_k^r\right)v_k
  =a(\eps_k)v_k.
\]
Векторы $v_k$ образуют базис, поэтому найден полный набор собственных
векторов. Формула для спектра также непосредственно следует из теоремы
об отображении спектра и равенства $\Sp(V)=\{\eps_0,\dots,\eps_{n-1}\}$.
\end{proof}

Обозначим через $D$ матрицу оператора $A$ в базисе собственных векторов:
\[
  D=\diag\bigl(a(\eps_0),a(\eps_1),\dots,a(\eps_{n-1})\bigr).
\]
Поскольку столбцы $F$ --- собственные векторы, имеем $AF=FD$. Следовательно,
\[
  \boxed{F^{-1}AF=D},\qquad
  \boxed{A=FDF^{-1}}.
\]
Таким образом, одна и та же матрица $F$ диагонализует все циклические
матрицы данного размера. Собственные числа $A$ могут совпадать; это не
мешает существованию базиса из векторов $v_k$.

\section{Критерий цикличности}

\begin{theorem}
Матрица $A\in\CC^{n\times n}$ является циклической тогда и только тогда,
когда она перестановочна с матрицей циклического сдвига:
\[
  \boxed{AV=VA}.
\]
\end{theorem}

\begin{proof}
Докажем обе импликации.

\smallskip
\noindent
\emph{Необходимость.} Если $A$ циклическая, то $A=a(V)$. Любая степень
$V$ перестановочна с $V$, поэтому
\begin{align*}
  AV&=(a_0E+a_1V+\dots+a_{n-1}V^{n-1})V\\
    &=a_0V+a_1V^2+\dots+a_{n-1}V^n\\
    &=V(a_0E+a_1V+\dots+a_{n-1}V^{n-1})=VA.
\end{align*}

\smallskip
\noindent
\emph{Достаточность.} Пусть $AV=VA$. Обозначим стандартные базисные
векторы через $e_0,\dots,e_{n-1}$, а столбцы $A$ --- через
$A^{(j)}=Ae_j$. Для $j=0,\dots,n-2$ выполнено $e_{j+1}=Ve_j$, откуда
\[
  A^{(j+1)}=Ae_{j+1}=AVe_j=VAe_j=VA^{(j)}.
\]
Значит, каждый следующий столбец получается сдвигом предыдущего.
Кроме того, $Ve_{n-1}=e_0$, поэтому $VA^{(n-1)}=A^{(0)}$.
Это и есть структура циклической матрицы.
\end{proof}

Смысл критерия состоит в том, что действие циклической матрицы согласовано
со сдвигом: сначала применить $A$, а затем сдвинуть результат --- то же
самое, что сначала сдвинуть входной вектор, а затем применить $A$.

\section{Алгебра циклических матриц}

\begin{proposition}
Множество циклических матриц размера $n\times n$ замкнуто относительно
сложения, умножения на скаляр и матричного умножения. Любые две такие
матрицы перестановочны. Если циклическая матрица обратима, её обратная
также циклическая.
\end{proposition}

\begin{proof}
Пусть $AV=VA$, $BV=VB$ и $\alpha\in\CC$. Тогда
\[
  (\alpha A)V=V(\alpha A),\qquad
  (A+B)V=V(A+B).
\]
Для произведения получаем
\[
  (AB)V=A(BV)=A(VB)=(AV)B=(VA)B=V(AB).
\]
По критерию цикличности матрицы $\alpha A$, $A+B$ и $AB$ циклические.

Для доказательства перестановочности воспользуемся символами:
\[
  A=a(V),\qquad B=b(V).
\]
Многочлены от одной матрицы перестановочны, поскольку любые её степени
перестановочны друг с другом. Следовательно,
\[
  AB=a(V)b(V)=b(V)a(V)=BA.
\]

Наконец, если $A$ обратима, умножим $AV=VA$ слева и справа на $A^{-1}$:
\[
  A^{-1}(AV)A^{-1}=A^{-1}(VA)A^{-1},
  \qquad VA^{-1}=A^{-1}V.
\]
Снова применяя критерий, получаем цикличность $A^{-1}$.
\end{proof}

\begin{definition}
Подалгебру матричной алгебры называют \emph{полной}, если вместе с каждой
своей матрицей, обратимой в полной матричной алгебре, она содержит и её
обратную.
\end{definition}

Итак, в терминологии лекции циклические матрицы образуют коммутативную
полную подалгебру алгебры всех матриц. Единичная матрица принадлежит этой
подалгебре. Здесь слово <<полная>> означает именно замкнутость относительно
взятия обратной, а не утверждение о метрической полноте.

\section{Обратимость и формула для обратной матрицы}

\begin{theorem}
Циклическая матрица $A=a(V)$ обратима тогда и только тогда, когда
\[
  a(\eps_k)\ne0\qquad\text{для всех }k=0,1,\dots,n-1.
\]
В этом случае
\[
  \boxed{A^{-1}=F\diag\left(
    \frac{1}{a(\eps_0)},\dots,\frac{1}{a(\eps_{n-1})}
  \right)F^{-1}}.
\]
\end{theorem}

\begin{proof}
Матрица $A=FDF^{-1}$ обратима тогда и только тогда, когда обратима
диагональная матрица $D$. Для этого необходимо и достаточно, чтобы все
её диагональные элементы были ненулевыми. Обратная к $D$ получается
заменой каждого диагонального элемента на обратный. Поэтому
\[
  (FDF^{-1})^{-1}=FD^{-1}F^{-1},
\]
что и даёт указанную формулу.
\end{proof}

Проверять символ требуется только в корнях $n$-й степени из единицы.
Его нули в других точках комплексной плоскости на обратимость $A$ не влияют.

\section{Арифметика индексов по модулю \texorpdfstring{$n$}{n}}

На множестве $\ZZ_n=\{0,1,\dots,n-1\}$ определим операции
\[
  k\oplus\ell=(k+\ell)\bmod n,\qquad
  k\odot\ell=(k\ell)\bmod n.
\]
Результатом всегда считаем представителя из $\{0,\dots,n-1\}$.
Например, при $n=5$
\[
  3\oplus3=1,\qquad 3\odot2=1.
\]
С этими операциями $\ZZ_n$ является коммутативным кольцом. Если $n$ ---
простое число, это поле. Для определения свёртки используется аддитивная
группа этого кольца.

Также введём обозначение для разности индексов:
\[
  k\ominus\ell=(k-\ell)\bmod n.
\]
Например, при $n=5$ имеем $1\ominus3=3$. При фиксированном $k$
отображение $\ell\mapsto k\ominus\ell$ переставляет элементы $\ZZ_n$.
Именно это позволяет менять переменную суммирования в формулах свёртки.

\section{Свёртка векторов}

\begin{definition}
\emph{Циклической свёрткой} векторов $a,b\in\CC^n$ называется вектор
$c=a*b\in\CC^n$, координаты которого определяются формулой
\[
  \boxed{(a*b)_k=\sum_{\ell=0}^{n-1}a_{k\ominus\ell}b_\ell},
  \qquad k=0,1,\dots,n-1.
\]
\end{definition}

Вычитание в индексе выполняется по модулю $n$. Поэтому, например, индекс
$-1$ означает $n-1$; выход за границы вектора не возникает.

\subsection{Коммутативность}

\begin{proposition}
Для любых $a,b\in\CC^n$ выполнено $a*b=b*a$.
\end{proposition}

\begin{proof}
В сумме для $k$-й координаты сделаем замену $p=k\ominus\ell$.
Тогда $\ell=k\ominus p$, и $p$ пробегает все элементы $\ZZ_n$:
\[
  (a*b)_k
  =\sum_{\ell=0}^{n-1}a_{k\ominus\ell}b_\ell
  =\sum_{p=0}^{n-1}a_pb_{k\ominus p}
  =\sum_{p=0}^{n-1}b_{k\ominus p}a_p
  =(b*a)_k.
\]
Равенство выполнено для каждой координаты, поэтому векторы совпадают.
\end{proof}

\subsection{Ассоциативность}

\begin{proposition}
Для любых $a,b,c\in\CC^n$ выполнено $(a*b)*c=a*(b*c)$.
\end{proposition}

\begin{proof}
Сначала раскроем левую часть:
\begin{align*}
  ((a*b)*c)_k
  &=\sum_{\ell=0}^{n-1}(a*b)_{k\ominus\ell}c_\ell\\
  &=\sum_{\ell=0}^{n-1}\sum_{m=0}^{n-1}
    a_{k\ominus\ell\ominus m}b_mc_\ell.
\end{align*}
Все суммы конечны, поэтому порядок суммирования можно менять.
В правой части получаем
\begin{align*}
  (a*(b*c))_k
  &=\sum_{r=0}^{n-1}a_{k\ominus r}(b*c)_r\\
  &=\sum_{r=0}^{n-1}\sum_{\ell=0}^{n-1}
    a_{k\ominus r}b_{r\ominus\ell}c_\ell.
\end{align*}
При фиксированном $\ell$ положим $m=r\ominus\ell$. Тогда
$r=m\oplus\ell$, а пара $(m,\ell)$ пробегает все пары индексов.
Следовательно,
\[
  (a*(b*c))_k
  =\sum_{\ell=0}^{n-1}\sum_{m=0}^{n-1}
    a_{k\ominus\ell\ominus m}b_mc_\ell
  =((a*b)*c)_k.
\]
\end{proof}

\subsection{Единичный элемент}

Положим $e=e_0=(1,0,\dots,0)^{\mathsf T}$. Тогда
\[
  (a*e)_k=\sum_{\ell=0}^{n-1}a_{k\ominus\ell}e_\ell
  =a_{k\ominus0}=a_k.
\]
В сумме остаётся только слагаемое с $\ell=0$. Значит,
\[
  a*e=e*a=a.
\]
Таким образом, $e_0$ является единицей относительно свёртки.

\subsection{Дистрибутивность и умножение на скаляр}

Для каждой координаты имеем
\begin{align*}
  (a*(b+c))_k
  &=\sum_{\ell=0}^{n-1}a_{k\ominus\ell}(b_\ell+c_\ell)\\
  &=\sum_{\ell=0}^{n-1}a_{k\ominus\ell}b_\ell
    +\sum_{\ell=0}^{n-1}a_{k\ominus\ell}c_\ell\\
  &=(a*b)_k+(a*c)_k.
\end{align*}
Следовательно,
\[
  a*(b+c)=a*b+a*c.
\]
Аналогично для $\alpha\in\CC$
\[
  (\alpha a)*b=\alpha(a*b)=a*(\alpha b).
\]
Итак, пространство $\CC^n$ с обычными сложением и умножением на скаляр
и с умножением, заданным свёрткой, образует коммутативную алгебру с единицей.

\section{Оператор свёртки}

\begin{definition}
При фиксированном $a\in\CC^n$ отображение
\[
  T_a\colon\CC^n\longrightarrow\CC^n,\qquad T_ax=a*x,
\]
называется \emph{оператором свёртки} с вектором $a$.
\end{definition}

Из дистрибутивности и согласованности с умножением на скаляр следует
линейность:
\[
  T_a(\alpha x+\beta y)=\alpha T_ax+\beta T_ay.
\]
В рукописи оператор обозначен буквой $A$; здесь обозначение $T_a$
позволяет отличать отображение от его матрицы.

\begin{remark}
Связь с циклическими матрицами можно увидеть непосредственно из формулы
для координат:
\[
  (T_ax)_k=\sum_{\ell=0}^{n-1}a_{k\ominus\ell}x_\ell.
\]
Значит, матрица $T_a$ в стандартном базисе имеет элементы
$A_{k\ell}=a_{k\ominus\ell}$. Это циклическая матрица с первым столбцом
$a$. Таким образом, $T_ax=Ax$, где $A=a(V)$. Это пояснение связывает два
раздела лекции.
\end{remark}

\section{Подчинённые матричные нормы}

В заключительной части рукописи рассматриваются нормы произвольной
матрицы $A=(a_{k\ell})\in\CC^{n\times n}$.

\begin{definition}
Матричная норма, подчинённая выбранной векторной норме, определяется как
\[
  \norm{A}=\sup_{x\ne0}\frac{\norm{Ax}}{\norm{x}}
  =\sup_{\norm{x}=1}\norm{Ax}.
\]
\end{definition}

Второе равенство следует из нормировки $x$: любой ненулевой вектор можно
заменить на $x/\norm{x}$. В конечномерном пространстве единичная сфера
компактна, поэтому этот супремум достигается.

\subsection{Норма, подчинённая сумме модулей координат}

Пусть $\norm{x}_1=\sum_{k=0}^{n-1}|x_k|$.

\begin{theorem}
\[
  \boxed{\norm{A}_1=\max_{0\leqslant\ell<n}
    \sum_{k=0}^{n-1}|a_{k\ell}|}.
\]
То есть $\norm{A}_1$ равна максимальной сумме модулей элементов столбца.
\end{theorem}

\begin{proof}
Обозначим правую часть через $M$. По неравенству треугольника
\begin{align*}
  \norm{Ax}_1
  &=\sum_{k=0}^{n-1}\left|\sum_{\ell=0}^{n-1}a_{k\ell}x_\ell\right|\\
  &\leqslant\sum_{\ell=0}^{n-1}
    \left(\sum_{k=0}^{n-1}|a_{k\ell}|\right)|x_\ell|\\
  &\leqslant M\sum_{\ell=0}^{n-1}|x_\ell|=M\norm{x}_1.
\end{align*}
Следовательно, $\norm{A}_1\leqslant M$.

Чтобы получить равенство, выберем столбец с номером $\ell_0$, на котором
достигается максимум, и положим $x=e_{\ell_0}$. Тогда $\norm{x}_1=1$,
а $Ax$ совпадает с этим столбцом. Поэтому
\[
  \norm{Ae_{\ell_0}}_1=\sum_{k=0}^{n-1}|a_{k\ell_0}|=M.
\]
Значит, $\norm{A}_1\geqslant M$, что завершает доказательство.
\end{proof}

\subsection{Норма, подчинённая максимуму модулей координат}

Для $\norm{x}_\infty=\max_k|x_k|$ справедлива формула
\[
  \boxed{\norm{A}_\infty=\max_{0\leqslant k<n}
    \sum_{\ell=0}^{n-1}|a_{k\ell}|}.
\]
Это максимальная сумма модулей элементов строки.

Для полноты поясним достижение этой оценки. Пусть $k_0$ --- номер строки
с максимальной суммой, а $a_{k_0\ell}=|a_{k_0\ell}|e^{i\theta_\ell}$
для её ненулевых элементов. Выберем
\[
  x_\ell=\begin{cases}
    e^{-i\theta_\ell},&a_{k_0\ell}\ne0,\\
    1,&a_{k_0\ell}=0.
  \end{cases}
\]
Тогда $\norm{x}_\infty=1$ и
\[
  (Ax)_{k_0}=\sum_{\ell=0}^{n-1}|a_{k_0\ell}|.
\]
Обратная оценка следует из неравенства треугольника для каждой строки.

\subsection{Евклидова норма и сингулярные числа}

Пусть $\norm{x}_2=(\sum_k|x_k|^2)^{1/2}$. Обозначим через
$A^*=\overline{A}^{\mathsf T}$ сопряжённо-транспонированную матрицу.
Матрица $A^*A$ эрмитова и неотрицательно определена, поскольку
\[
  x^*A^*Ax=\norm{Ax}_2^2\geqslant0.
\]
Если её собственные числа равны $\lambda_0,\dots,\lambda_{n-1}$, то числа
\[
  \sigma_k=\sqrt{\lambda_k},\qquad k=0,\dots,n-1,
\]
называются сингулярными числами $A$. Все они неотрицательны.
По формуле максимума отношения Рэлея
\[
  \norm{A}_2^2
  =\max_{x\ne0}\frac{x^*A^*Ax}{x^*x}
  =\lambda_{\max}(A^*A).
\]
Следовательно,
\[
  \boxed{\norm{A}_2=\max_k\sigma_k
    =\sqrt{\lambda_{\max}(A^*A)}}.
\]

\subsection{Связь трёх норм}

В рукописи приведена оценка
\[
  \boxed{\norm{A}_2\leqslant
    \sqrt{\norm{A}_1\norm{A}_\infty}}.
\]
Дополним запись её доказательством. По неравенству Коши--Буняковского
для каждой строки
\begin{align*}
  \left|\sum_{\ell=0}^{n-1}a_{k\ell}x_\ell\right|^2
  &\leqslant\left(\sum_{\ell=0}^{n-1}|a_{k\ell}||x_\ell|\right)^2\\
  &\leqslant\left(\sum_{\ell=0}^{n-1}|a_{k\ell}|\right)
    \left(\sum_{\ell=0}^{n-1}|a_{k\ell}||x_\ell|^2\right).
\end{align*}
Суммируя по $k$, получаем
\begin{align*}
  \norm{Ax}_2^2
  &\leqslant\norm{A}_\infty
    \sum_{\ell=0}^{n-1}\left(\sum_{k=0}^{n-1}|a_{k\ell}|\right)|x_\ell|^2\\
  &\leqslant\norm{A}_\infty\norm{A}_1\norm{x}_2^2.
\end{align*}
Делим на $\norm{x}_2^2$ при $x\ne0$, берём супремум и извлекаем корень.

\section{Заключительные замечания и темы докладов}

\subsection{Число обусловленности}

В рукописи указан заголовок <<Число обусловленности>>, но дальнейшее
изложение этой темы отсутствует. Для связи с предыдущей лекцией напомним
только определение: для обратимой матрицы $A$
\[
  \cond(A)=\norm{A}\,\norm{A^{-1}}.
\]
Обе нормы берутся в одном и том же соглашении. Это справочное дополнение
к записи, а не восстановленная по скану часть доказательства.

\subsection{Диагонализация и сингулярное разложение}

В конце записи рядом с $A=FDF^{-1}$ есть краткое замечание о типе
разложения. Поясним различие. Формула $A=FDF^{-1}$ выражает
\emph{диагонализацию}: диагональные элементы $D$ --- собственные числа,
которые могут быть комплексными. В сингулярном разложении
\[
  A=U\Sigma W^*
\]
матрицы $U,W$ унитарны, а диагональные элементы $\Sigma$ --- вещественные
неотрицательные сингулярные числа. Поэтому запись $A=FDF^{-1}$ сама по
себе не является записью сингулярного разложения.

\subsection{Темы докладов}

\begin{enumerate}
  \item Дискретное преобразование Фурье: свойства матрицы $F$ с
        доказательствами.
  \item Быстрое преобразование Фурье.
\end{enumerate}

\end{document}
